% F3: respostas estruturadas reais (final_results.jsonl). Trechos longos abreviados com [...] e "...".
% Usa \ttfamily e \char34 (e não verbatim) porque o babel torna a aspa dupla ativa e o verbatim quebra acentos.
\providecommand{\jq}{\char34{}}
\providecommand{\ji}{\hspace*{1.2em}}
\vspace{2mm}
\noindent\begin{minipage}{\columnwidth}
\ttfamily\scriptsize\frenchspacing
\renewcommand{\arraystretch}{1.0}
\begin{tabular}{@{}l@{}}
\textrm{\itshape a02, \armTextToCypher{}} \\
\{ \\
\ji\jq rationale\jq: \jq Contei [...] Adrien Rabiot \\
\ji\ji lidera com 10, seguido de [...]\jq, \\
\ji\jq players\jq: [\jq Adrien Rabiot\jq], \\
\ji\jq value\jq: 10.0, \\
\ji\jq no\_data\jq: false \\
\} \\[1.5ex]
\textrm{\itshape a02, \armGraphTools{}} \\
\{ \\
\ji\jq rationale\jq: \jq Consultei [...] Adrien Rabiot \\
\ji\ji liderou com 10 passes [...]\jq, \\
\ji\jq players\jq: [\jq Adrien Rabiot\jq, \\
\ji\ji \jq Kylian Mbappé Lottin\jq, ...], \\
\ji\jq value\jq: 10.0, \\
\ji\jq no\_data\jq: false \\
\} \\[1.5ex]
\textrm{\itshape u02 (controle), \armTextToCypher{}} \\
\{ \\
\ji\jq rationale\jq: \jq O grafo não contém \\
\ji\ji distância percorrida [...]\jq, \\
\ji\jq players\jq: [], \\
\ji\jq value\jq: null, \\
\ji\jq no\_data\jq: true \\
\}
\end{tabular}
\captionof{figure}{Respostas estruturadas reais do \modelHaiku{}. Os trechos abreviados aparecem como [...] e as reticências indicam nomes omitidos.}\label{fig:json}
\end{minipage}\par\smallskip
